% 0_abstract.tex
Multi-agent discussion systems let several role-conditioned agents
deliberate on a question, and their evaluation increasingly relies on a
single large language model (LLM) acting as judge: the judge reads the
discussion transcript and scores quality dimensions, and the resulting
scores drive claims about which system components matter. The robustness
of such claims to the identity of the judge has not been examined. We
audit this evaluation protocol on \sysname{}, a decentralized roundtable
framework in which each agent bids for the floor with a self-assessed
confidence score. We re-implemented the framework's protocol verbatim as
a headless, dependency-free testbed, generated 33 discussions under two
driver models, and scored every transcript with two judges from
different model families. Three findings emerge. First, the two judges
produce near-zero rank agreement on discussion quality (Spearman
$\rho=-0.05$; per-dimension agreement is weak and negative on insight
depth), despite both coming from the same vendor---suggesting the
disagreement is deeper than vendor-specific bias. Second, the dominant
inter-judge effect is a systematic leniency offset (bootstrap 95\%
confidence intervals excluding zero) rather than self-preference, which
we do not detect; this refines the standard advice that using a
different judge family suffices. Third, component-importance conclusions
drawn from a single-judge ablation reverse direction under the other
judge (Cliff's $\delta$ up to $-1.0$), and in a same-machine
replication where only the judge varies, neither judge ranks the full
system above all ablations---the audited evaluation's component
ordering is not reproduced under any judge--driver combination tested.
We conclude that single-judge LLM scores do not support
component-importance claims about multi-agent discussion systems in
their current form, and we propose a minimum multi-judge reporting
norm. All code, transcripts, judgments, and an independent jury record
are released with the paper.
